\documentclass[11pt]{article}
\usepackage{fontspec}
\setmainfont{Latin Modern Roman}
\usepackage{amsmath,amssymb,mathtools,unicode-math}
\setmathfont{Latin Modern Math}
\usepackage{xurl,hyperref}
\hypersetup{hidelinks,pdftitle={Problem statement}}
\newcommand{\sref}[2]{\hyperlink{source:#1:#2}{[S#2]}}
\newcommand{\cataloguescope}[1]{\paragraph{Mathematical question.}#1}
\begin{document}
\section{Research Problems in Function Theory — Problem 5.19}\label{statement}

\subsection{Definitions and mathematical statement}

\paragraph{Shared notation.}$\mathbb N=\{1,2,\ldots\}$, and $\mathbb Z,\mathbb Q,\mathbb R,\mathbb C$ denote the integers, rationals, reals and complex numbers. Any other conventions are specified locally.


Write $\mathbb D=\{z\in\mathbb C:|z|<1\}$ and $\mathbb T=\partial\mathbb D$. All Taylor coefficients below are taken at $0$.A meromorphic function $f$ on $\mathbb D$ has an asymptotic value $a\in\widehat{\mathbb C}=\mathbb C\cup\{\infty\}$ if $f(\gamma(t))\to a$ in the spherical metric along a continuous path $\gamma:[0,1)\to\mathbb D$ with $|\gamma(t)|\to1$. The path is not required to converge to a single boundary point. A value $a$ is locally exceptional at $\zeta\in\mathbb T$ if, in some neighborhood of $\zeta$ in $\mathbb D$, it is assumed only finitely many times; after shrinking the neighborhood it is omitted there.
\paragraph{Statement recorded in the supplied source.}
Does there exist a meromorphic function $f$ on $\mathbb D$ with no asymptotic values, but with a locally exceptional value at some $\zeta\in\mathbb T$? \sref{statement}{2}
\subsection{Short English statement}
Historical question: can a meromorphic disc function without any asymptotic value still omit a value near one boundary point? The source update reports examples.
\subsection{Sources}
\begin{description}
\item[\hypertarget{source:statement:1}{S1}] ULAM.AI and the UnsolvedMath Contributors, UnsolvedMath: A Curated Collection of Open Mathematics Problems, catalogue entry. CC BY 4.0. \url{https://huggingface.co/datasets/ulamai/UnsolvedMath}.
\item[\hypertarget{source:statement:2}{S2}] W. K. Hayman and E. F. Lingham, Research Problems in Function Theory (2018), arXiv:1809.07200, Chapter 5, Problem 5.19, its update and the preceding shared notation. \url{https://arxiv.org/abs/1809.07200}.
\end{description}
\label{end:statement}
\end{document}
